\section{进一步的问题：允许拒绝是否足以消除定位歧义？}
\label{sec:rejection-question}

应用接口增加了按半径拒绝的功能，但“可以拒绝”本身不是可靠性的证明。本节在命题\ref{prop:mirror}的两世界构造上推导一个精确限制，再把后续研究问题写成可检验的约束。这是对既有不可区分性论证的直接延伸，不主张文献首次性；它属于报告中的数学分析，不计入智能体工具已经执行的数学检查次数。

\begin{proposition}[镜像歧义下的接受率与错误率限制]
\label{prop:reject-mirror}
设两个可能世界的当前位置为$x_0,x_1$，距离$D=\|x_0-x_1\|>0$，且完整观测$Z$的分布完全相同。两个世界使用同一可测规则，以$Z$及独立于$Z$和世界的同分布内部随机性$U$为输入，输出接受指标$A\in\{0,1\}$和位置$\widehat x$（拒绝时可任意定义该位置）。令共同接受概率为$p=\Pr(A=1)>0$。对任意$0\le r<D/2$，有
\begin{align}
 \max_{j\in\{0,1\}}\Pr_j\{\|\widehat x-x_j\|>r\mid A=1\}&\ge\frac12,\label{eq:reject-conditional}\\
 \max_{j\in\{0,1\}}\Pr_j\{A=1,\ \|\widehat x-x_j\|>r\}&\ge\frac p2.\label{eq:reject-joint}
\end{align}
因此，若要求每个世界的“接受且错误”概率均不超过$\alpha<1/2$，必须满足$p\le2\alpha$，即拒绝率至少为$1-2\alpha$。这里拒绝既不记为成功，也不记为一次已接受输出的错误；式\eqref{eq:reject-joint}采用的正是这个联合事件口径。
\end{proposition}

\begin{proof}
在两个世界使用相同分布的$(Z,U)$，其中$U$是估计器随机性。由于$2r<D$，两个闭球$B(x_0,r)$与$B(x_1,r)$互不相交。因此，对每次接受的相同输出，两个世界中至多有一个满足误差不超过$r$，即
\[
 \mathbf1_{\{A=1,\|\widehat x-x_0\|>r\}}+
 \mathbf1_{\{A=1,\|\widehat x-x_1\|>r\}}
 \ge \mathbf1_{\{A=1\}}.
\]
对共同分布积分，两个联合错误概率之和至少为$p$，其最大值至少为$p/2$。再除以两个世界共同的正接受概率$p$，得到条件结论。最后由联合错误均不超过$\alpha$推出$p\le2\alpha$。
\end{proof}

上述$1/2$界在允许规则知道两个候选点的两点模型中可以达到：以概率$p$接受，接受时等概率返回$x_0$或$x_1$。对任一世界，其条件错误率恰为$1/2$，联合错误率恰为$p/2$。这只是证明常数无法普遍收紧的理论规则，不是部署建议；它也不依赖声称某个局部最小值更可信。当$p=0$时联合错误为零，但接受后的条件错误率没有定义，不能将全拒绝解释为条件准确率100\%。

该结论说明三件事。第一，无信息的歧义不能仅通过缩小接受集合变成较高的最坏世界条件准确率；两个世界始终看到相同的接受事件。第二，允许大量拒绝可以降低联合错误，却会降低服务率。例如$\alpha=0.1$时，必要拒绝率至少为80\%。第三，这不否定前文的共形边际保证：保证所采用的分布、随机性、区域宽度及拒绝计分方式不同，不能互相替换。若输出区域允许同时覆盖两个镜像点，就必须承受相应的几何宽度。

\subsection{可供后续研究的受约束风险问题}
要提出有内容的改善目标，需要先排除上述不可区分性，或明确其出现概率。这里的$P$应定义为完整观测$Z$与真实当前位置、运动参数的\emph{联合分布}，而不是只给观测分布一个场景名称。固定有限观测时间窗，令$\mathcal P_0$中的轨迹以概率一为$x(t)=x$，$\mathcal P_1$中的轨迹以概率一为$x(t)=x+tv$且$v_{\min}\le\|v\|\le v_{\max}$，其中$0<v_{\min}\le v_{\max}<\infty$。位置参数域、污染机制、上报误差、观测数与时间设计必须在比较前固定；以下风险还要求相应条件二阶矩有限。

一种可进一步明确的\emph{全局}几何约束是：给定允许的理想锚点--时间设计$\mathcal D$，在规定参数域上，无噪声距离映射$F_{\mathcal D}(x,v)=(\|x+t_i v-a_i\|)_i$满足
\[
 \|F_{\mathcal D}(x,v)-F_{\mathcal D}(x',v')\|_2
 \ge c\|x-x'\|_2,\qquad c>0.
\]
它比较参数域内所有允许的参数对，因而排除了位置不同但理想距离完全相同的镜像；单独限制某拟合点Jacobian的最小奇异值不能替代它。哪些设计满足该条件、噪声与污染后能保留怎样的可辨识度，仍需证明；本报告不声称现有合成场景已经满足这一条件，也不把它本身当作有噪声风险界。

对规则$(\widehat x,A)$定义接受率及接受后的均方风险
\[
 a_P=\Pr_P(A=1),\qquad
 R_P=\mathbb E_P[\|\widehat x-x\|^2\mid A=1].
\]
接受指标仍只能依赖观测与预先指定的内部随机性。先固定基线估计器；其风险$R_P^{\rm WLS\mid A}$与$R_P^{\rm ref\mid A}$也必须在\emph{同一个由待评规则决定的接受事件}上计算；不能让各方法各自挑选不同的容易案例后直接比较。设$C$为统一口径下单次定位与接受判定的计算成本。一个明确的后续问题是：在
\[
 \inf_{P\in\mathcal P_0\cup\mathcal P_1}a_P\ge a_0>0,
 \qquad
 \sup_{P\in\mathcal P_0\cup\mathcal P_1}\mathbb E_P[C]\le B
\]
时，是否存在规则同时满足
\begin{align}
 \sup_{P\in\mathcal P_0}(R_P-R_P^{\rm WLS\mid A})&\le\delta_0,\\
 \sup_{P\in\mathcal P_1}(R_P-R_P^{\rm ref\mid A})&\le-\delta_1,
 \qquad \delta_1>0?
\end{align}
问题中的$(\delta_0,\delta_1,a_0,B)$给出静态容忍、移动收益、服务率和计算预算四项取舍。精度使用均方误差，因而不能把本报告的均值误差比较直接当作这些约束的证明。共同接受事件只保证比较同一批案例；若进一步声称等计算量优势，还须对基线施加同口径的预算。不同参数下可行集合可能为空：例如若移动分布类包含参照方法在接受事件上风险为零的分布，统一的正收益$\delta_1>0$立即不可能。命题\ref{prop:reject-mirror}提供另一类由不可区分性导致的障碍。

本次实验展示了促成该后续问题的具体失败机制，并实现了测量这些取舍的工具；它没有求出上述可行域，也没有给出统一的最优拒绝规则。下一步应先在明确的二维几何与污染分布类中求下界，再设计加入方向信息或轨迹先验的规则，最后进行预算匹配的数值检验。
